\section{Background and summary}\label{intro}

The Rub' al-Khali (Empty Quarter) spans approximately 500,000 km$^2$
of the southern Arabian Peninsula \cite{mandaville_1986}. Its soils
experience large temperature fluctuations and irregular rainfall
\cite{nelli2024wind}, yet support diverse microbial communities
\cite{makhalanyane2015}. Repeated observations across soil compartments
and seasons provide a basis for studying these communities together
with their environmental conditions.

Previous Saudi Arabian surveys examined elevation gradients
\cite{yasir2015saudi}, rainfall responses across soil compartments
\cite{aslam2016rainfall}, and bulk and rhizosphere soils across regions
\cite{khan2020saudi}. Studies of desert plants further connected
microbial profiles with host-associated habitats \cite{eida2018plants}.
A repeated transect adds a specific data-management requirement:
each community profile must remain connected to its field specimen,
collection visit, laboratory processing and associated measurements.
Site identity alone is insufficient for temporal comparisons because
a site can have several visits and assays can use different field
replicates.

Multi-campaign resources such as Tara Oceans link campaign, station
and sampling-event registries to environmental and sequence data
\cite{taraoceans2015}. MIxS/MIMARKS specify contextual fields for
environmental sequence data \cite{mixs}, and STREAMS extends reporting
guidance for environmental microbiomes \cite{streams2025}. The National
Microbiome Data Collaborative uses persistent identifiers and shared
schemas to connect environmental microbiome data \cite{nmdc2026}.
Knowledge graphs such as KG-Microbe \cite{kgmicrobe2026} and
MetagenomicKG \cite{metagenomickg2026} organize taxonomic, functional and
environmental information across sources. We apply these principles
to the specimen and process provenance of a single landscape survey.

We collected soil and plant material during five campaigns
(2023--2025) along a 1,043~km transect across the Saudi Rub' al-Khali.
The catalogue contains 60 primary sites and 10 named locations.
A 2,550-row source ledger contains 2,516 non-control specimen rows.
The canonical amplicon table contains 1,271 profiles; the companion
ecology analysis retains 1,237 after control and read-depth exclusions.
Associated data comprise 274 field environmental logs, 712 admitted
archived-soil pH observations, 71 field XRF sessions, 725 laboratory
XRF records and monthly climate summaries. These counts describe
different observational units, whose links and coverage are documented
in the source ledgers.

We encoded sites, specimens, processes, qualities and measurement
values with the Semanticscience Integrated Ontology (SIO) \cite{sio}.
ChEBI \cite{chebi}, ENVO \cite{envo} and NCBI Taxonomy
\cite{ncbitaxon} provide chemical cross-references, environmental
classes and validated taxonomic identifiers. Original classifier
lineages and unresolved mappings retain their source provenance.
The resulting RDF modules support queries across these linked
observations.

Technical validation examines source reconciliation, measurement
admission, identifier mappings, structural constraints and executable
query examples. We report the scope of each check separately,
including the constructs and modules processed by ELK.
The portal at \url{https://rubalkhali.science/} exposes the graph
through site pages and a SPARQL endpoint. Data Records and Data
Availability identify the distributed artifacts and outstanding
release dependencies. The companion ecology manuscript reports the
community, environmental-association and functional analyses
\cite{ecology_companion}; this Data Descriptor documents their source
data, representation and reuse conditions.
